% citum-example.tex -- Feature demonstration of the Citum LuaLaTeX package
%
% NOTE: This is an entirely artificial document created for feature illustration.
% All arguments, references, and authors are fabricated. No scholarly claims are made.
%
% Usage:
%   ../scripts/citum-lualatex citum-example.tex
%
% The package is pipe-only: LuaLaTeX records a format_document JSON-RPC request
% per pass. The compile script runs citum-server externally by default; pass
% --shell-escape to let LuaLaTeX spawn citum-server directly.
%
% For multilingual entries (Japanese, Spanish) full Unicode rendering requires
% appropriate fonts. The data model is correct regardless; install e.g. Noto fonts
% and set them via \setmainfont / \newfontfamily as needed.

\documentclass{article}
\usepackage[margin=1.2in]{geometry}
\usepackage{fontspec}
\usepackage{hyperref}

% Load the Citum citation package.
%   style   = path to a Citum YAML style file (see demo/ for apa-7th.yaml)
%   bibfile = path to a Citum YAML bibliography
%   locale  = optional BCP 47 locale tag
%   server  = optional explicit path to citum-server binary (pipe transport)
\usepackage[
    style   = apa-7th,
    bibfile = example-refs,
    % locale  = en-US,
    % server  = /usr/local/bin/citum-server,
]{citum}

\title{The Infrastructure of Scholarly Memory\\
       {\large A Citum Lua\LaTeX{} Feature Demonstration}}
\author{Citum}
\date{2026}

\begin{document}

\maketitle

% -----------------------------------------------------------------------
\section{Introduction: citation as institution}
% -----------------------------------------------------------------------

Scholarly citation serves as the primary mechanism through which academic
knowledge is organized, attributed, and transmitted.
As \textcite{chen2017} has argued, the practices surrounding citation
reflect deeper epistemological commitments about what counts as knowledge
and who produces it.
The formal apparatus of bibliographic reference --- the footnote, the
parenthetical, the numbered list --- constitutes what \textcite{webb1968}
called ``the footnote as institution,'' embedding social relations within
the seemingly neutral mechanics of scholarly communication.

% -----------------------------------------------------------------------
\section{Non-integral citations with locators}
% -----------------------------------------------------------------------

Non-integral (parenthetical) citations are produced with \verb|\cite|.
A bare key places author and year in parentheses \cite{chen2017}.
Optional locators follow the usual prefix convention:

\begin{itemize}
  \item Page: \cite[p.~42]{chen2017}
  \item Chapter: \cite[ch.~3]{webb1968}
  \item Section: \cite[sect.~2]{okafor2024}
\end{itemize}

Multiple keys in a single parenthetical use \verb|\cites|:
\cites{chen2017, okafor2024}.

% -----------------------------------------------------------------------
\section{Integral (narrative) citations}
% -----------------------------------------------------------------------

Integral citations are produced with \verb|\textcite|, which places the
author name grammatically in the sentence and appends the year in parentheses.

\textcite{okafor2024} provide large-scale empirical support for this view,
showing through computational analysis that topological centrality in citation
networks correlates more strongly with institutional affiliation than with
measured scholarly impact.

\textcite[ch.~2]{chen2020} extends the analysis to the digital transition,
arguing that platform architectures reproduce existing hierarchies of
credibility.
% Note: same author as chen2017. In author-date styles both entries are
% disambiguated by year. In note-based styles, name memory abbreviates the
% second citation to the short form.

% -----------------------------------------------------------------------
\section{Archival and non-standard sources}
% -----------------------------------------------------------------------

Citum's \texttt{archive-info} field models collection hierarchy and shelfmark.
The date \texttt{"1891?"} is an EDTF Level~1 uncertain date, which renders
as ``1891?'' in APA and as ``um~1891'' in a German locale.

\textcite{harrington1891} observed --- in manuscript notes preserved at the
British Library --- that bibliographic method was already contested among
Victorian scholars of natural history \cite[ff.~12--15]{harrington1891}.

% -----------------------------------------------------------------------
\section{Multilingual sources and EDTF approximate dates}
% -----------------------------------------------------------------------

The production of scholarly citation apparatus is not confined to Anglophone
contexts.
\textcite{tanaka_yuki2019} demonstrate that citation practices in Japanese
science studies reflect distinct epistemological conventions.
Parallel observations emerge from Spanish-language scholarship
\cite{garcia2022}.
% garcia2022 uses EDTF "2022~" (approximate date), which renders as
% "ca.~2022" in en-US and as "vers~2022" in fr-FR.

% -----------------------------------------------------------------------
\section{Complex grouped citations}
% -----------------------------------------------------------------------

For citations with per-item locators or prefixes, use the standard multi-citation
commands:

\cites[see]{chen2017}[also]{okafor2024}

Integral grouped citations use \verb|\textcites|:

\textcites{chen2017}{chen2020}
have both contributed to this line of inquiry.

Alternatively, you can also use the legacy \verb|\citestart| / \verb|\citeitem| / \verb|\citeend| block format:

\citestart
  \citeitem[see]{chen2017}
  \citeitem[also]{okafor2024}
\citeend

% -----------------------------------------------------------------------
\section{Bibliography}
% -----------------------------------------------------------------------

% Bibliography is split by entry type, mirroring the HTML demo:
%   type=manuscript    → primary sources (harrington1891)
%   not-type=manuscript → secondary sources (all others)

\subsection*{Primary Sources}
\addcontentsline{toc}{subsection}{Primary Sources}
\printcitumbibliography[type=manuscript]

\subsection*{Secondary Sources}
\addcontentsline{toc}{subsection}{Secondary Sources}
\printcitumbibliography[not-type=manuscript]

\end{document}
